% Dua belas faktor Twelve-Factor App dikelompokkan menjadi tiga kelompok,
% masing-masing dengan satu pertanyaan bisnis yang mudah diingat.
\begin{tikzpicture}[
  font=\scriptsize,
  execute at begin node={\hyphenpenalty=10000},
  induk/.style={draw, rounded corners, align=center, minimum height=0.8cm,
                text width=4.2cm, font=\scriptsize\bfseries},
  kelompok/.style={draw, rounded corners, align=left, text width=3.1cm,
                   minimum height=2.7cm, anchor=north, inner sep=4pt},
  tanya/.style={align=center, text width=3.3cm, font=\tiny\itshape,
                text=gray!70!black, anchor=north},
  panah/.style={-{Stealth[length=2mm]}, thick}
]
  \node[induk, fill=praditagreen!15] (induk) at (0,1.6)
    {Twelve-Factor App\\{\tiny\mdseries 12 faktor, 3 kelompok}};

  \node[kelompok, fill=blue!8] (a) at (-3.9,0)
    {\textbf{A. Kode, konfigurasi, dan \textit{backing service}}\\[3pt]
     I \textit{Codebase}\\ II \textit{Dependencies}\\ III \textit{Config}\\
     IV \textit{Backing services}};
  \node[kelompok, fill=orange!12] (b) at (0,0)
    {\textbf{B. \textit{Build}, proses, dan skala}\\[3pt]
     V \textit{Build, release, run}\\ VI \textit{Processes}\\
     VII \textit{Port binding}\\ VIII \textit{Concurrency}\\ IX \textit{Disposability}};
  \node[kelompok, fill=blue!8] (c) at (3.9,0)
    {\textbf{C. Lingkungan, \textit{log}, dan tugas administratif}\\[3pt]
     X \textit{Dev/prod parity}\\ XI \textit{Logs}\\ XII \textit{Admin processes}};

  \node[tanya] at ([yshift=-3pt]a.south)
    {Dapatkah satu kode dipasang di banyak lokasi tanpa diubah?};
  \node[tanya] at ([yshift=-3pt]b.south)
    {Dapatkah kapasitas ditambah dan dikurangi dalam hitungan menit?};
  \node[tanya] at ([yshift=-3pt]c.south)
    {Dapatkah sistem dirawat, diaudit, dan dipulihkan dengan cepat?};

  \draw[panah] (induk.south) -- ++(0,-0.3) -| (a.north);
  \draw[panah] (induk.south) -- (b.north);
  \draw[panah] (induk.south) -- ++(0,-0.3) -| (c.north);
\end{tikzpicture}
